% future/HTMtableRCU.tex
% SPDX-License-Identifier: CC-BY-SA-3.0

\begin{table*}
\centering
\small
\input{future/HTMtableColor}
\setstretch{0.95}
\setlength{\tabcolsep}{4pt}\OneColumnHSpace{-.9in}
\ebresizewidth{
\resizebox{6.5in}{!}{
\begin{tabularx}{6.8in}{p{0.95in}cXcX}
\toprule
  &
    & \multicolumn{1}{c}{用户空间RCU或Hazard Pointer增强的锁} &
      & \multicolumn{1}{c}{Hardware Transactional Memory} \\
\midrule
  基本思想 &
    & 一次只允许一个线程访问给定的一组对象。 &
      & 使对给定对象集的操作以原子方式执行。 \\
\midrule
  适用范围 &
    & \Pl 处理所有操作。 &
      & \Pl 处理可撤销操作。 \\
\addlinespace[4pt]
  &
    & &
      & \Mn 不可撤销操作迫使回退（通常回退到锁）。 \\
\midrule
  可组合性 &
    & \Pl 读取端仅受宽限期等待操作限制。 &
      & \Dw 受不可撤销操作、事务大小和死锁限制
        （假设基于锁的回退代码）。\\
\addlinespace[4pt]
  &
    & \Mn 更新端受死锁限制。读取端减少死锁。 &
      & \\
\midrule
  可扩展性与性能 &
    & \Mn 数据必须可分区以避免更新端之间的锁竞争。 &
      & \Mn 数据必须可分区以避免冲突。 \\
\addlinespace[4pt]
  &
    & \Pl 读取端不需要分区。 &
      & \\
\cmidrule{3-5}
  &
    & \Dw 更新端的分区通常必须在设计时固定。 &
      & \Pl 自动动态调整分区，精确到缓存行边界。 \\
\cmidrule{3-5}
  &
    & \Pl 读取端不需要分区。 &
      & \Mn 回退需要分区（对于罕见的回退不太重要）。 \\
\cmidrule{3-5}
  &
    & \Dw 更新端锁原语通常导致昂贵的缓存未命中
      和内存屏障指令。 &
      & \Mn 事务开始/结束指令通常不会导致缓存未命中，
        但确实有内存排序和开销影响。 \\
\cmidrule{3-5}
  &
    & \Pl 更新端竞争效应集中在获取和释放上，
      因此临界区以全速运行。 &
      & \Mn 竞争会中止冲突的事务，即使它们已经
        运行了很长时间。 \\
\addlinespace[4pt]
  &
    & \Pl 读取端不与更新端或彼此竞争。 &
      & \\
\cmidrule{3-5}
  &
    & \Pl 读取端原语通常是有界无等待的，开销低。
      （Hazard pointer为无锁，开销低。） &
      & \Mn 只读事务受冲突和回滚影响。
        除回退代码提供的之外没有前向进展保证。 \\
\cmidrule{3-5}
  &
    & \Pl 当数据仅对更新端可见时，私有化操作
      简单、直观、高效且可扩展。 &
      & \Mn 私有化数据增加事务大小。 \\
\addlinespace[4pt]
  &
    & \Mn 对于读取端可见的数据，私有化操作开销较大
      （但仍然直观且可扩展）。 &
      & \\
\midrule
  硬件支持 &
    & \Pl 通用硬件即可。 &
      & \Mn 需要新硬件（正开始变得可用）。 \\
\cmidrule{3-5}
  &
    & \Pl 性能不受缓存几何细节影响。 &
      & \Mn 性能严重依赖缓存几何结构。 \\
\midrule
  软件支持 &
    & \Pl API已存在，有大量代码和经验，调试器正常工作。 &
      & \Mn API正在出现，DBMS之外经验很少，事务中的
        断点可能有问题。 \\
\midrule
  与其他机制的交互 &
    & \Pl 有长期成功交互的经验。 &
      & \Dw 刚刚开始研究交互。 \\
\midrule
  实际应用 &
    & \Pl 有。 &
      & \Pl 有。 \\
\midrule
  广泛适用性 &
    & \Pl 有。 &
      & \Mn 尚无定论。 \\
\bottomrule
\end{tabularx}
}}
\caption{锁（由RCU或Hazard Pointer增强）与HTM的比较
  （\colorbox{plus}{优势}、\colorbox{minus}{劣势}、
  \colorbox{down}{严重劣势}）}
\label{tab:future:Comparison of Locking (Augmented by RCU or Hazard Pointers) and HTM}
\end{table*}
